\section{Arrival Geometry and Minimax Risk}\label{sec:geometry}
\subsection{Model and Information}
Arms have labels in a countable set. Fix a deterministic sequence of integer arrival times $\tau_i\ge1$, with
\[
 A_t=\{i:\tau_i\le t\},\qquad 1\le K_t:=|A_t|<\infty.
\]
At round $t$, the learner observes the entire new batch before choosing $I_t\in A_t$. Arm $i$ supplies a potential reward stream $(X_{i,n})_{n\ge1}$ of independent and identically distributed random variables with values in $[0,1]$ and mean $\mu_i$. Streams are independent across arms. The $n$th pull reveals $X_{i,n}$; all other rewards remain unobserved. Policy randomization is independent of these streams. The choice at a round may depend only on the revealed arrival sets, previous actions and observations, and this randomization.

Let $B_t=\max_{i\in A_t}\mu_i$ and $\Delta_i(t)=B_t-\mu_i$. We use the random pseudo-regret
\begin{equation}\label{eq:regret}
 R_T=\sum_{t=1}^T\Delta_{I_t}(t),\qquad 0\le R_T\le T.
\end{equation}
Both $B_t$ and the gap of an already available arm are nondecreasing. No further assumption on these quantities is imposed. In particular, a new batch can contain several optimal arms or change the benchmark by any amount compatible with the reward range.

For a fixed schedule $\tau$ and deterministic endpoint $T$, define
\begin{equation}\label{eq:minimax}
 \cR_T(\tau)=\inf_{\pi\in\Pi_{T,\tau}}\sup_P\EE_P^\pi R_T.
\end{equation}
The class $\Pi_{T,\tau}$ permits knowledge of $T$ and of the complete arrival schedule, but not of the reward laws. The supremum ranges over the independent bounded stationary laws just described. Giving the lower-bound learner this additional information strengthens the comparison. Our upper-bound policy belongs to the smaller class that uses no future arrival information and no endpoint.

The minimax value is unchanged if the supremum is restricted to Bernoulli rewards. For the nontrivial direction, replace each observed bounded reward $X$ by $\ind\{U\le X\}$ with an independent uniform random variable $U$. The resulting stream for arm $i$ consists of independent Bernoulli variables with mean $\mu_i$, and the regret of the induced actions depends on the same means. Taking an arbitrarily near-optimal Bernoulli policy proves the assertion. This reduction concerns minimax risk; it does not identify the instance-specific information in a general reward law.

\subsection{The Lifetime Spectrum}
Choose any arm $a\in A_1$ as an anchor for counting. Its mean is unknown and it need not be optimal. Define
\[
 H_i=T-\tau_i+1,\qquad L=\lceil\log_2T\rceil,\qquad b_\ell=2^\ell.
\]
Among the arms other than $a$, let $G_0=\{i:H_i=1\}$ and
\[
 G_\ell=\{i:2^{\ell-1}<H_i\le2^\ell\},\qquad 1\le\ell\le L,
 \qquad k_\ell=|G_\ell|.
\]
The scale costs and their sum are
\begin{equation}\label{eq:spectrum}
 s_{\ell,T}=\min\{b_\ell,\sqrt{k_\ell b_\ell}\},\qquad
 \cS_T=\sum_{\ell=0}^{L}s_{\ell,T}.
\end{equation}
Every possible anchor has lifetime $T$, so the spectrum is independent of the anchor's identity. Empty groups contribute zero; if $K_T=1$, the spectrum and regret are both zero.

\begin{theorem}[Schedule-specific minimax risk]\label{thm:main}
There is one causal policy $\pi^{\mathrm{MC}}$, specified in Section~\ref{sec:policy}, such that for every allowed infinite arrival schedule and every deterministic $T\ge1$,
\begin{equation}\label{eq:main}
 \frac{\cS_T}{10000}\le\cR_T(\tau)
 \le\sup_P\EE_P R_T(\pi^{\mathrm{MC}})
 \le\min\{T,300\cS_T\}.
\end{equation}
The policy does not use $T$, future arrivals, the eventual number of arms or reward means. Its sample counts and reward sums are never erased.
\end{theorem}

The constants are not optimized. Throughout, a rate characterization or the notation $\asymp$ refers to multiplicative bounds by positive universal constants, unless stated otherwise. It does not refer to an exact asymptotic minimax constant.

The truncation in $s_{\ell,T}$ matters when a batch contains more arms than the number of remaining rounds. Such a batch cannot contribute more regret than its remaining lifetime. When its size is smaller, its cost has the familiar square-root form. The theorem states that these costs accumulate across scales, despite the fact that the sampling windows overlap. Section~\ref{sec:policy} gives the upper-bound mechanism, and the supplement proves the compatible lower construction.

\subsection{A Capacity Without a Grid}
Write $M=K_T-1$. A fractional allocation to the competing arms is feasible precisely when
\begin{equation}\label{eq:occupancy}
 \cN_T=\left\{n\in[0,\infty)^M:
 \sum_{i:H_i\le h}n_i\le h\ \text{for }h=1,\ldots,T\right\}.
\end{equation}
Indeed, the arms in the displayed sum are available only in the last $h$ time slots. Conversely, process arms from shortest to longest lifetime and place their demands in the unused part of their legal suffixes. The nested constraints ensure that the available capacity is sufficient at each step. Unused capacity can be assigned to the anchor.

Define
\begin{equation}\label{eq:capacity}
 \phi(x)=\min\{x,\sqrt{x}\},\qquad
 \cC_T=\max_{n\in\cN_T}\sum_{i\ne a}\phi(n_i),
\end{equation}
with $\cC_T=0$ when $M=0$. The feasible set is compact and $\phi$ is continuous, so a maximizer exists. The capacity is computed from arrival times alone. This calculation does not establish a regret guarantee for a policy based on the maximizing allocation.

\begin{theorem}[Capacity characterization]\label{thm:capacity}
For every schedule and endpoint,
\begin{equation}\label{eq:capacity-equivalence}
 \tfrac14\cS_T\le\cC_T\le\cS_T,
 \qquad
 \frac{\cC_T}{10000}\le\cR_T(\tau)
 \le\sup_P\EE_P R_T(\pi^{\mathrm{MC}})
 \le\min\{T,1200\cC_T\}.
\end{equation}
\end{theorem}

To see the upper comparison, a group has total allocated mass at most $b_\ell$, and its objective is at most both that mass and the square root of $k_\ell$ times that mass. For the other direction, allocate $b_\ell/(4k_\ell)$ to each arm of a nonempty group. The geometric sum of group capacities verifies every suffix constraint, and each group contributes at least $s_{\ell,T}/4$. The complete argument is in the supplement.

The linear part of $\phi$ prevents a fractional observation from contributing more than its mass. Replacing $\phi$ by an untruncated square root would assign unbounded capacity to a sufficiently large batch arriving in the last round. In contrast, \eqref{eq:capacity} gives exactly one for every nonempty such batch.
